\documentclass[11pt,a4paper,sans]{moderncv}
\moderncvstyle{classic}
\moderncvcolor{blue}
\usepackage[utf8]{inputenc}
\usepackage[scale=0.85]{geometry}
\nopagenumbers

% ---------------------------------------------------------------------------
% VARIANTE DIRECIONADA — Trilho B: AI Infrastructure / Observability / Evals
% Base: curriculum.tex (EN conciso). Mesmos empregos, mesmas datas, mesmos fatos.
% Mudanca de ENFASE: observabilidade GenAI e "own the output" sobem para o topo.
% Nao substitui curriculum.tex — e uma variante adicional.
% ---------------------------------------------------------------------------

\name{Edimar Calebe}{Castanho}
\title{Cloud \& Platform Engineer — AI Infrastructure, Observability \& Evals}
\address{Curitiba, Brazil}{Open to relocation (Europe \& US)}
\email{contato@calebe.dev.br}
\phone[mobile]{+55~(41)~99827~1302}
\social[linkedin]{calebe94}
\social[github]{calebe94}
\extrainfo{\faRss\href{https://blog.calebe.dev.br/}{blog.calebe.dev.br}}

\begin{document}
\maketitle

\section{Summary}
\cvitem{}{
Cloud, platform and backend engineer with 7+ years building production infrastructure and the observability around it.
Delivered GenAI observability end-to-end as a contributor to OpenTelemetry v2 span normalization for LLM, agent, tool and workflow traces, alongside a unified model cost and pricing override system serving multi-tenant APIs.
Leads a vSphere-to-AWS migration covering zero-trust networking, least-privilege IAM, Packer AMI automation and ephemeral CI runners, with a \textasciitilde\$566 quarterly cost projection against a \$15k annual cap.
Directs AI coding agents across infrastructure and application work, reviewing generated output for correctness, security and test coverage before merge.
Proficient in AWS, OpenStack and GCP with Terraform, Ansible and CI/CD automation (GitHub Actions, GitLab CI).
Bachelor of Engineering in Computer Engineering. Fluent in English and Portuguese, conversational Spanish.
}

\section{Professional Experience}

\cventry{May 2026 - Present}{Senior Cloud Engineer (Contract)}{Illumina (via Insight Global)}{Remote}{}{\begin{itemize}
\item Directing AI agents to scaffold and iterate on Terraform/AWS infrastructure, reviewing generated code for correctness, security posture and adherence to canonical baselines before merge — owning the output rather than the generation
\item Leading zero-trust vSphere-to-AWS migration for the Instrumentation Software team, provisioning GitHub Enterprise ephemeral auto-scaling runners and Terraform-driven E2E test environments (multi-VPC, multi-instance, complex network rules) within a \$15k annual budget cap
\item Architected AWS infrastructure blueprint with subnet segmentation, least-privilege IAM roles, SSM-managed instance profiles, and Packer-based AMI automation (Rocky Linux 9), with cost modeling projecting \textasciitilde\$566 Q3 spend
\item Delivered proof-of-concept ephemeral runner pool with single-ENI topology, budget allocation tagging, pre-warm scheduling, and governance guardrails (CODEOWNERS review enforcement, yamllint CI gates, automated smoke-test workflows)
\item Coordinated cross-team architecture alignment with DevOps and Platform Solutions engineers, validating pre-wired IAM roles, security groups, and AMI IDs against canonical infrastructure baselines before provisioning
\end{itemize}}

\cventry{Dec 2025 - May 2026}{Senior Software Engineer (Contract)}{Galileo.ai (acquired by Cisco), via Eversynced}{Remote}{}{\begin{itemize}
\item Developed Mastra AI framework extension for OpenTelemetry v2, normalizing GenAI spans and enabling distributed tracing across agent, tool and workflow execution — the observability layer that makes production LLM systems debuggable and auditable
\item Extended the Galileo Python SDK with OpenTelemetry span serialization for LLM, agent, tool, and workflow span types, expanding distributed tracing coverage across the platform
\item Consolidated model cost calculation into a single source of truth, migrating pricing data to a cached cost table and adding per-organization price overrides to the scoring pipeline
\item Built end-to-end multi-tenant model pricing override system with row-level security (FastAPI, SQLAlchemy, Alembic), backed by a test-driven DAO layer and enforcing tenant isolation across three platform services
\item Implemented custom integration features across the full stack, including model aliasing with legacy backward compatibility, encrypted HTTP headers, and API key authentication
\item Directed AI coding agents to implement features across the stack, reviewing their output for correctness, architecture, and test coverage before merge
\end{itemize}}

\cventry{Jan 2023 - Nov 2025}{Senior Cloud \& Backend Engineer}{LuizaLabs (Magalu Cloud)}{Curitiba, Brazil}{}{\begin{itemize}
\item Developed asynchronous backend services (asyncio, FastAPI, Pydantic) exposing RESTful APIs to automate OpenStack VM lifecycle management, reducing operational overhead 30\%
\item Led Python infrastructure automation for VM provisioning on OpenStack (Glance, Nova, Cinder), Incus/LXD, and KVM/QEMU, improving platform reliability across 50K+ instances
\item Automated Ubuntu and Windows Server image builds with Ansible and QCOW2 compression, cutting deployment time 40\%
\item Designed Terraform modules for performance engineering and benchmarking environments (sysbench, fio, iperf), deployed via GitLab CI/CD and GitHub Actions
\item Maintained PostgreSQL- and MySQL-backed microservices, collaborating with DevOps teams on code quality and infrastructure scaling
\end{itemize}}

\cventry{Mar 2020 - Dec 2022}{Embedded \& Backend Systems Analyst}{T2I Group}{Curitiba, Brazil}{}{\begin{itemize}
\item Architected IoT data pipeline using Flask, MQTT, and FreeRTOS, aggregating real-time sensor data from 200+ roadside units
\item Designed RESTful APIs for vehicle detection on ESP32 modules and TM4C microcontrollers, processing 5K+ daily classifications
\item Consolidated 50+ repositories onto GitLab CI/CD and GitHub Actions pipelines, reducing build times 30\%
\end{itemize}}

\cventry{Feb 2019 - Feb 2020}{Backend Developer Intern}{Alta Rail Technology}{Colombo, Brazil}{}{\begin{itemize}
\item Developed real-time train monitoring dashboard with Django and Elasticsearch for telemetry analytics, and optimized satellite communication, reducing transmission costs 25\%
\end{itemize}}

\cventry{Jan 2017 - Dec 2017}{Embedded Systems Intern}{T2I Group}{Curitiba, Brazil}{}{\begin{itemize}
\item Built OBD-II fleet tracking device (PIC18F, GSM/GPRS) with a web-based monitoring backend and Raspberry Pi automation system
\end{itemize}}

\section{Technical Skills}
\cvitem{AI \& Observability}{OpenTelemetry v2 GenAI span normalization (LLM, Agent, Tool, Workflow), distributed tracing, model cost and pricing override engines, agent orchestration and output verification, multitenant model-serving APIs, Observability-Driven Development}
\cvitem{Languages}{Python, TypeScript, JavaScript, SQL, C/C++, Bash}
\cvitem{Backend}{FastAPI, Django, Flask, Express.js, SQLAlchemy, Pydantic, Celery, asyncio, pytest, PostgreSQL, MySQL, MongoDB, Elasticsearch, SQLite, Redis, REST API Design, Microservices}
\cvitem{Cloud \& DevOps}{AWS (EC2, S3, Lambda, IAM, VPC, SSM), OpenStack (Glance, Nova, Cinder), GCP, KVM/QEMU, Incus, Packer, Terraform, Ansible, Docker, Kubernetes, GitLab CI/CD, GitHub Actions, GitHub Enterprise, ArgoCD, Infrastructure as Code, Zero-Trust Security, Cloud Migration}
\cvitem{Spoken Languages}{Portuguese (Native), English (Fluent), Spanish (Conversational)}

\section{Education}
\cventry{2015 - 2023}{Bachelor of Engineering in Computer Engineering}{Universidade Positivo}{Curitiba, Brazil}{}{}
\cventry{2011 - 2013}{Electromechanical Technical Course}{Instituto Federal do Paran\'a}{Curitiba, Brazil}{}{}

\end{document}
